\newcommand{\GeometryStatus}{Exact helix signature and generic period-one helper checked.}
\newcommand{\PellStatus}{Exact original signature checked; no new helper.}
\newcommand{\MomentsStatus}{Symbolic design; arithmetic endpoints checked in Python.}
\newcommand{\RowsStatus}{Source-reviewed recognition design; Lean integration unverified.}
\newcommand{\ValidationSummary}{Two retained standalone Lean files completed successfully: the exact helical theorem with a generic period-one helper, and the exact Pell threshold theorem. The seam adaptation is supplied as a source patch supported by the checked generic helper; its full production module was not rebuilt. The symbolic signal-moment adapter timed out without diagnostics and is explicitly unvalidated. The canonical-row and edge-recognition files are proposed integrations, not a completed Lean check. An optional full production dependency rebuild was stopped before completion; no whole-project build, runtime improvement or structural net line saving is claimed.}
\newcommand{\GeometryResults}{The exact-signature \code{HelicalPeriod.lean} compiles with Lean 4.32.0 and pinned Mathlib. The two proposed production theorem blocks change from 19 to 8 lines and from 8 to 4 lines. The complete two-file candidate changes from 253 to 238 physical lines, a 15-line reduction, with no new production helper. Of these 15 lines, 11 belong to the exact checked helix replacement and 4 to the seam adaptation whose production integration remains unchecked. These are source-patch measurements, not a full-module build certificate.}
\newcommand{\PellResults}{The exact-signature \code{PellAnchorSmall.lean} compiled successfully (exit 0), importing only \code{Mathlib.NumberTheory.PellMatiyasevic}. It neither imports the project theorem being replaced nor adds a helper declaration. The production patch changes the declaration from 15 to 11 physical lines and its proof body from 13 to 9 lines: 4 lines removed. The patch applies cleanly to the pinned checkout. The full \code{NativePellBootstrap} module was not rebuilt. The axiom printout is \code{propext}, \code{Classical.choice}, \code{Quot.sound}; no stronger axiom-preservation claim is made.}
\newcommand{\MomentsResults}{The included Python computation reproduces all four existing signal-family endpoints, and the formulas were checked algebraically. The proposed \code{SignalMomentsCore.lean} compiler attempt timed out after 300 seconds without diagnostics; this is an unvalidated design artifact, not a proof success or a demonstrated proof error. The rail-rank sum, initial-segment adapter, and complete four-client migration remain implementation work. No achieved source saving or Lean performance improvement is claimed.}
\newcommand{\RowsResults}{The supplied \code{CanonicalOperatingRow.lean} retains the complete receiver equality, states the interior canonical-row theorem, derives the south-color theorem by the construction API, and gives an immediate row-uniqueness corollary. The dependency plan has been source-reviewed, but the file was not compiled against a rebuilt production import graph. The canonical theorem still partitions old head, moved head and other positions; moving that partition is not by itself a reduction. The stronger API and reuse of color semantics are the benefit to test. Charge the additional statements and any compatibility wrappers before claiming net shortening.}
\newcommand{\ReproductionNotes}{Use \code{scripts/verify.py} in a checkout with the pinned Lean/Mathlib dependencies. Its default mode checks the two completed standalone files and verifies patch applicability without applying patches. The optional proposal mode also tries the unvalidated signal, canonical-row and edge-recognition files after their production imports have been built. The packet records exact inputs and outcomes; local cache locations from this review are not reproduction requirements.}
